\section*{Chapter \ref{ch:b3:development-microeconomics} --- Development Microeconomics}

\begin{solution}{exo:b3:development-microeconomics:1}
$6.81/94.81 = 7.2~\%$ of the control mean. Temptation goods, $-0.79$ dollars per two weeks. % figdata: microcredit
\end{solution}

\begin{solution}{exo:b3:development-microeconomics:2}
The poor are 2, 4 and 6, with shortfalls 0.75, 0.5 and 0.25 of the line. $P_0 = 3/5 = 0.6$; $P_1 = 1.5/5 = 0.3$; $P_2 = (0.5625 + 0.25 + 0.0625)/5 = 0.175$.
\end{solution}

\begin{solution}{exo:b3:development-microeconomics:3}
$6 + 4 + 2 = 12$, which is $N z P_1 = 5 \times 8 \times 0.3$.
\end{solution}

\begin{solution}{exo:b3:development-microeconomics:4}
Effects on assets and income proportional to the size of the transfer, with no
jump: households given small transfers keeping the same share of them as those given
large ones, and no size after which households keep accumulating.
\end{solution}

\begin{solution}{exo:b3:development-microeconomics:5}
The tested solution is \code{ex05_fgt_kenya.py}, in the chapter's \code{exercises/}
folder. At 3.00 dollars: 45.0, 13.6 and 5.7~\%; at 4.20: 65.0, 25.9 and 12.8~\%; at 8.30: 90.0, 53.0 and 34.9~\%. % output: ex05_fgt_kenya
\end{solution}

\begin{solution}{exo:b3:development-microeconomics:6}
Weights $1/16$, $1/4$, $1/16$, summing to 0.375: $(10/16 + 2/4 + 6/16)/0.375 = 1.5/0.375 = 4.0$, with standard error $\sqrt{1/0.375} = 1.63$.
\end{solution}

\begin{solution}{exo:b3:development-microeconomics:7}
Without \omterm{def:b3:development-microeconomics:microcredit}{group liability} both loans are repaid with probability $0.8^2 = 0.64$ and each with 0.8; with it, each
is repaid unless both fail, $1 - 0.2^2 = 0.96$, provided $2r \le 3$. Break-even: $r = 1/0.8 = 1.25$ without, $r = 1/0.96 = 1.04$ with.
\end{solution}

\begin{solution}{exo:b3:development-microeconomics:8}
Per 100 children the programme pays $80 \times 400 = 32{,}000$ a year for 10 extra child-years: 3{,}200 dollars each, because 70 of the 80
paid would have enrolled anyway.
\end{solution}

\begin{solution}{exo:b3:development-microeconomics:9}
Prices of local goods rising in treated villages relative to controls, physical
quantities consumed (food, durables) unchanged, and non-recipients in treated
villages consuming less than controls.
\end{solution}

\begin{solution}{exo:b3:development-microeconomics:10}
Transfer sensitivity is proved in \cref{prop:b3:development-microeconomics:fgt}. For $P_0$, take $y_i = 1$, $y_j = z - 0.1$ and $t = 0.2$: after the
transfer $j$ is above the line, so $P_0$ falls by $1/N$, although the poorest person is poorer.
\end{solution}

\begin{solution}{exo:b3:development-microeconomics:11}
The tested solution is \code{ex11_heterogeneity.py}, in the chapter's
\code{exercises/} folder. Fixed effect 6.31 (standard error 1.72); $Q = 12.87$ on 6 degrees of freedom; $\tau^2 = 25.62$ ($\tau = 5.06$);
random effects 6.89 with interval $[1.57, 12.21]$, wider than the fixed-effect one. % output: ex11_heterogeneity
\end{solution}

\begin{solution}{exo:b3:development-microeconomics:12}
$P^{1 + \varepsilon} = 1.15$: prices rise 15~\% with $\varepsilon = 0$ and $1.15^{0.1} - 1 = 1.4~\%$ with $\varepsilon = 9$. The Kenyan trial found minimal inflation, as with elastic supply.
\end{solution}

\begin{solution}{pb:b3:development-microeconomics:cash}
% figdata: cash_stats
\textbf{1.} 7{,}925 dalens a year; the line is 4{,}755.

\textbf{2.} $P_0 = 15.8~\%$, $P_1 = 2.54~\%$ and $P_2 = 0.57~\%$ of the 4{,}899 people surveyed.

\textbf{3.} 592{,}450 dalens a year, which is $N z P_1$.

\textbf{4.} 195.

\textbf{5.} Consumption varies less from year to year than income and records what
households actually live on, including what they draw from savings.

\textbf{6.} It rewards helping those just below the line and ignores the depth of
poverty that $P_1$ and $P_2$ measure.

\textbf{7.} 97 treated and 98 controls.

\textbf{8.} 3{,}878 against 3{,}927 dalens: random assignment makes the groups alike before the transfer, up to
chance.

\textbf{9.} 4{,}371 and 4{,}092.

\textbf{10.} 280 dalens per person, standard error 88.

\textbf{11.} $[107, 453]$.

\textbf{12.} 29.9~\% of the treated and 11.2~\% of the controls.

\textbf{13.} Consumption grows on average, varies from year to year, and exceeds income
for many poor households, so some controls end above the line.

\textbf{14.} Controls 1{,}919; transfers of 3{,}000: 3{,}327; of 6{,}000: 9{,}867; of 9{,}000: 14{,}238 dalens.

\textbf{15.} 8.3~\%, 93.5~\% and 100~\%.

\textbf{16.} With 3{,}000, most households keep only 55~\% of assets plus transfer; with 9{,}000, all grow them by 25~\%.

\textbf{17.} Yes: the share of assets kept jumps between the 3{,}000 and 6{,}000 transfers, at the 7{,}000 threshold.
Proportional effects would keep the same share at every size.

\textbf{18.} It makes crossing the threshold depend on a random draw, not on households'
ability or plans.

\textbf{19.} 11.1, 29.0 and 53.3~\%, against 11.2~\% for controls.

\textbf{20.} Transfers too small to cross the threshold are largely spent and leave no lasting gain;
fewer, larger transfers lift more households for the same budget.

\textbf{21.} $1.1 \times 564{,}000 = 620{,}400$ dalens.

\textbf{22.} $(0.299 - 0.112) \times 97 = 18.1$.

\textbf{23.} $620{,}400/18.1 = 34{,}253$ dalens.

\textbf{24.} $3{,}000/0.08 = 37{,}500$ dalens.

\textbf{25.} The training effect comes from another study and population; lifting above
the line ignores how far below it households were and how long the gain lasts; and
18 households is a small and noisy count.

\textbf{26.} After two years the transfer raised consumption by 280 dalens per person (95~\% interval 107 to 453), at a
cost of 34{,}253 dalens per household lifted above the line.
\end{solution}
